\documentclass[10pt]{article}
\usepackage[T1]{fontenc}
\usepackage[utf8]{inputenc}
\usepackage[french]{babel}
\usepackage{amsmath,amssymb,mathrsfs,mathpazo}
\usepackage{geometry}
\geometry{a4paper,margin=2cm}
\usepackage{graphicx,array,tabularx,tabularray,colortbl,multirow,multicol,booktabs,enumitem,pifont,ulem,fancybox,diagbox,cancel}
\usepackage{pstricks,pstricks-add,pst-plot,pst-tree,pst-node,pst-text,pst-func,pst-math,pst-bspline,pst-3dplot,pst-eucl,pst-solides3d}
\let\pscalculate\undefined
\usepackage{pst-bezier}
\usepackage{tikz,pgfplots,tkz-tab}
\usetikzlibrary{arrows,arrows.meta,patterns,calc,positioning,trees,decorations.markings,shapes.geometric,shapes.misc,shapes.arrows,intersections,decorations.pathmorphing,decorations.pathreplacing,backgrounds,shadows,babel}
\pgfplotsset{compat=1.18}
\usepgfplotslibrary{fillbetween}
\usepackage[np]{numprint}
\usepackage{listings,makecell,pas-tableur,twoopt}
\let\barre\undefined
\usepackage[pstricks]{tablvar}
\definecolor{darkgreen}{RGB}{0,100,0}
\providecommand{\red}{\color{red}}
\providecommand{\blue}{\color{blue}}
\providecommand{\green}{\color{green}}
\providecommand{\black}{\color{black}}
\providecommand{\white}{\color{white}}
\providecommand{\gray}{\color{gray}}
\providecommand{\cyan}{\color{cyan}}
\providecommand{\magenta}{\color{magenta}}
\providecommand{\yellow}{\color{yellow}}
\providecommand{\violet}{\color{violet}}
\providecommand{\orange}{\color{orange}}
\providecommand{\euro}{\texteuro}
\providecommand{\EUR}{\texteuro}
\providecommand{\ieme}{\textsuperscript{e}}
\providecommand{\degres}{\ensuremath{^\circ}}
\providecommand{\degree}{\ensuremath{^\circ}}
\providecommand{\vv}[1]{\overrightarrow{#1}}
\providecommand{\Coord}[1]{#1}
\providecommand{\decosix}{}
\providecommand{\starredbullet}{\ensuremath{\bullet}}
\pagestyle{empty}
\setlength{\parindent}{0pt}
\usepackage{adjustbox}\DeclareUnicodeCharacter{2236}{:}\DeclareUnicodeCharacter{2190}{\ensuremath{\leftarrow}}
\providecommand{\R}{\mathbb{R}}
\providecommand{\N}{\mathbb{N}}
\providecommand{\Z}{\mathbb{Z}}
\providecommand{\Q}{\mathbb{Q}}
\providecommand{\D}{\mathbb{D}}
\providecommand{\C}{\mathbb{C}}
\providecommand{\e}{{\mathrm{e}}}
\providecommand{\vect}[1]{\overrightarrow{#1}}
\providecommand{\vectt}[1]{\overrightarrow{\mathrm{#1}}}
\providecommand{\vvt}[1]{\overrightarrow{\mathrm{#1}}}
\providecommand{\np}[2][]{{#2~#1}}
\providecommand{\vv}[1]{{\overrightarrow{#1}}}
\renewcommand{\d}{{\mathrm{d}}}
\providecommand{\degres}{{^\circ}}
\providecommand{\degre}{{^\circ}}
\providecommand{\blue}{\color{blue}}
\providecommand{\red}{\color{red}}
\providecommand{\enskip}{\quad}
\providecommand{\llbracket}{\lbrack\!\lbrack}
\providecommand{\rrbracket}{\rbrack\!\rbrack}
\providecommand{\euro}{{\text{€}}}
\providecommand{\no}{{\text{n°}}}
\providecommand{\interval}[2]{{[#1;#2]}}
\providecommand{\intervFO}[2]{{[#1;#2[}}
\providecommand{\dotfill}{\ldots}

\providecommand{\ds}{\displaystyle}
\providecommand{\textding}[1]{#1}
\providecommand{\Oij}{\(\mathrm{O};\vec{\imath},\vec{\jmath}\)}
\providecommand{\Oijk}{\(\mathrm{O};\vec{\imath},\vec{\jmath},\vec{k}\)}
\providecommand{\Ouv}{\(\mathrm{O};\vec{u},\vec{v}\)}
\providecommand{\intervFF}[2]{\ensuremath{[#1;#2]}}
\renewcommand{\euro}{\ifmmode\text{\texteuro}\else\texteuro\fi}
\DeclareUnicodeCharacter{00B0}{\ifmmode^{\circ}\else\textdegree\fi}
\DeclareUnicodeCharacter{20AC}{\ifmmode\text{\texteuro}\else\texteuro\fi}
% Énoncés : Éducation nationale. Transcription et corrigés : APMEP, Denis Vergès et contributeurs. Découpage et classement : CoopMaths, Sébastien Lozano, Éric Elter et contributeurs. Adaptation de la présentation : Maths Explorer.
% CC BY-SA : https://coopmaths.fr/www/about/
% https://coopmaths.fr/alea/static/e3c/2020/tex/e3c_2020_00_sujet61_3.tex
\begin{document}
\noindent\textit{APMEP / CoopMaths. CC BY-SA. Présentation adaptée par Maths Explorer.}\par\medskip
\medskip

Un restaurant propose à sa carte deux types de dessert : un assortiment de macarons et une part de tarte tatin. Des études statistiques montrent que:

\setlength\parindent{1cm}
\begin{itemize}
\item l'assortiment de macarons est choisi par 50\,\% des clients;
\item la part de tarte tatin, est choisie par 30\,\% des clients;
\item 20\,\% des clients ne prennent pas de dessert;
\item aucun client ne prend plusieurs desserts.
\end{itemize}
\setlength\parindent{0cm}

Le restaurateur a remarqué que:

\setlength\parindent{1cm}
\begin{itemize}
\item parmi les clients ayant pris un assortiment de macarons, 80\,\% prennent un café;
\item parmi les clients ayant pris une part de tarte tatin, 60\,\% prennent un café;
\item parmi les clients n'ayant pas pris de dessert, 90\,\% prennent un café.
\end{itemize}
\setlength\parindent{0cm}

On interroge au hasard un client de ce restaurant. 

On note les évènements suivants:

\setlength\parindent{1cm}
\begin{itemize}
\item $M$ : \og Le client prend un assortiment de macarons \fg{};
\item $T$ : \og Le client prend une part de tarte tatin \fg{} ;
\item $P$ :\og Le client ne prend pas de dessert\fg{} ;
\item $C$ : \og Le client prend un café\fg{} et $\overline{C}$ l'évènement contraire de $C$.
\end{itemize}
\setlength\parindent{0cm}

\medskip

\begin{enumerate}
\item En utilisant les données de l'énoncé, préciser la valeur de $P(T)$ probabilité de $T$ et celle de $P_T(C)$ probabilité de l'évènement $C$ sachant que $T$ est réalisé.
\item  Recopier et compléter l'arbre ci-dessous :

\begin{center}
\pstree[treemode=R,nodesepB=3pt,levelsep=2.5cm]{\TR{}}
{\pstree{\TR{$M$~~}\taput{0,5}}
	{\TR{$C$} \taput{}
	\TR{$\overline{C}$} \tbput{}
	}
\pstree{\TR{$T$~~}\taput{}}
	{\TR{$C$} \taput{}
	\TR{$\overline{C}$} \tbput{}
	}
\pstree{\TR{$P$~~}\tbput{}}
	{\TR{$C$} \taput{}
	\TR{$\overline{C}$} \tbput{}
	}
}
\end{center}

\item  
	\begin{enumerate}
		\item Exprimer par une phrase ce que représente l'évènement $M \cap C$ puis calculer $P(M \cap C)$.
		\item Montrer que $P(C) = 0,76$.
	\end{enumerate}
\item  Quelle est la probabilité que le client prenne un assortiment de macarons sachant qu'il prend un café ? (On donnera le résultat arrondi au centième).
\end{enumerate}

\bigskip

\end{document}
